problem:  
Let \((X,d,\mathfrak{m})\) be an \(\mathrm{RCD}(0,N)\) space such that its isometry group \(G = \mathrm{Iso}(X,d)\) acts with **cohomogeneity one**, i.e. the orbits of the action are hypersurfaces of codimension one in \(X\). Assume all orbits are closed and that \(G\) acts effectively and isometrically with compact isotropy groups.  

**Question:**  
Is it true that such an \((X,d,\mathfrak{m})\) must be globally a smooth Riemannian manifold endowed with a \(G\)-invariant metric and measure (hence a classical cohomogeneity‑one manifold with nonnegative Ricci curvature), or can one construct a genuinely **non‑smooth cohomogeneity‑one \(\mathrm{RCD}(0,N)\) space**?  

In particular, does the synthetic \(\mathrm{RCD}(0,N)\) condition, when combined with one‑codimensional symmetry, already exclude the appearance of geometric singularities along the orbit space, or could singular “warped measure cones’’ arise that still satisfy \(\mathrm{RCD}(0,N)\)?  

Why is it a "good" problem:  
This problem moves beyond the fully homogeneous case—where classification is essentially understood—and attacks the next natural tier of symmetry: cohomogeneity one actions. These are geometrically simple yet deep enough to permit singularities (e.g., cone‑type or quotient singularities) in the synthetic setting. The question probes whether the analytic regularity implied by the \(\mathrm{RCD}\) condition still forces smooth bundle structures under almost‑maximal symmetry, or whether new non‑smooth models with group actions can exist.  

Resolving it would illuminate how lower Ricci curvature bounds interact with topological and symmetry‑induced stratifications, bridging geometric analysis, metric measure theory, and transformation group theory. It is simple to state yet profound in scope—an exemplary “narrow entrance, deep interior’’ problem that could clarify the precise border between synthetic and smooth homogeneous geometries.